\documentclass[12pt,a4paper]{article}
\usepackage[utf8]{inputenc}
\usepackage[T1]{fontenc}
\usepackage[spanish]{babel}
\usepackage{geometry}
\usepackage{longtable}
\usepackage{enumitem}
\usepackage{hyperref}
\usepackage{tabularx}
\usepackage{graphicx}
\geometry{margin=2.5cm}

\title{Sistema Nozama - Casos de Uso y Requisitos}
\author{Bryan Miguel Amador Castillo \\ Edson Joel Carrera Avila \\ Saul Lopez Mendoza}
\date{10/04/2025}

\begin{document}
\maketitle
\tableofcontents
\newpage

% ============================
% PLANTEAMIENTO DEL PROYECTO
% ============================

\section{Planteamiento}
Actualmente, muchos pequeños y medianos negocios comerciales no cuentan con
una plataforma digital que les permita ofrecer sus productos en línea de forma
organizada y eficiente. Los clientes se ven obligados a realizar compras
presenciales o mediante canales informales (WhatsApp, redes sociales) que no
garantizan la disponibilidad real del producto, lo que genera desconfianza, pedidos
duplicados y pérdida de ventas por falta de control de inventario.

Se requiere un sistema web que permita a los clientes explorar un catálogo de
productos organizados por categorías, y a los clientes registrados gestionar un
carrito de compras, realizar pagos y, de forma automática, actualizar el stock
disponible tras cada transacción. De igual manera, el negocio necesita un panel de
administración para gestionar productos (con captura por código UPC), categorías,
usuarios y órdenes de compra.

El sistema Nozama busca resolver la falta de plataformas confiables para pequeños
negocios mediante un e-commerce con gestión de inventario y administración centralizada.

% ============================
% ALCANCE DEL PROYECTO
% ============================

\section{Alcance del proyecto}
\subsection{Funcionalidades}
\begin{itemize}[label=--]
    \item Navegación pública del catálogo para usuarios no registrados (Invitados).
    \item Registro e inicio de sesión de clientes y administradores.
    \item Visualización del catálogo de productos por categorías.
    \item Búsqueda y filtrado de productos.
    \item Administración del carrito de compras (agregar, eliminar, modificar cantidad).
    \item Proceso de pago con resumen de orden.
    \item Actualización automática del stock tras cada compra confirmada.
    \item Panel de administración: gestión de productos con captura en formato UPC, categorías (CRUD) y usuarios.
    \item Historial de órdenes de compra por cliente.
    \item Control de diferentes tipos de clientes (estándar / VIP con descuentos).
    \item Manejo de promociones y descuentos por categoría o producto.
\end{itemize}

\subsection{Funcionalidades excluidas}
\begin{itemize}[label=--]
    \item Integración con pasarelas de pago reales (Stripe, PayPal, MercadoPago).
    \item Uso de la dirección del cliente.
    \item Aplicación móvil nativa.
    \item Envíos y logística (seguimiento de paquetes).
    \item Sistema de reseñas y calificaciones de productos.
    \item Chat en vivo con soporte.
\end{itemize}

\subsection{Límites técnicos}
\begin{itemize}[label=--]
    \item Backend: Python 3.x con Flask.
    \item Frontend: HTML5, CSS3, Bootstrap 5, Jinja2 Templates.
    \item Base de datos: MySQL con SQLAlchemy ORM.
    \item Autenticación: Flask-Login con sesiones.
    \item Sistema operativo de desarrollo: Ubuntu/Linux.
    \item Control de versiones: Git / GitHub.
\end{itemize}

\subsection{Límites de tiempo}
\begin{itemize}[label=--]
    \item Avance 1 (10/04/2026): Planteamiento, análisis de requisitos, casos de uso y mockups.
    \item Avance 2 (fecha por definir): Módulo de autenticación, catálogo y carrito implementados.
    \item Entrega final: Sistema completo con pruebas aplicadas (partición de equivalencia, valores límite, pruebas funcionales).
\end{itemize}

% ============================
% OBJETIVOS DEL PROYECTO
% ============================

\section{Objetivos}

\subsection{Objetivo General}
Desarrollar un sistema web de comercio electrónico con Python y Flask que permita
a los clientes explorar un catálogo de productos, gestionar un carrito de compras y
completar el proceso de pago, con actualización automática del inventario y un
panel de administración para la gestión del negocio.

\subsection{Objetivos específicos}
\begin{enumerate}[label=OE\arabic*:]
    \item Implementar un módulo de administración de usuarios con roles diferenciados (invitado, cliente estándar, cliente VIP, administrador).
    \item Desarrollar el catálogo de productos con organización por categorías y funcionalidad de búsqueda/filtrado accesible para usuarios no registrados.
    \item Implementar el carrito de compras con persistencia de sesión y cálculo automático de totales y descuentos.
    \item Desarrollar el flujo de pago que registre la orden y actualice el stock de cada producto vendido.
    \item Construir el panel de administración para la gestión de productos con captura en formato UPC, categorías y órdenes (altas, bajas, cambios).
    \item Aplicar técnicas de pruebas de software (partición equivalente, análisis de valores límite) sobre los módulos críticos del sistema.
\end{itemize}

\section{Análisis de requisitos}

% ============================
% ACTORES DEL SISTEMA
% ============================

\subsection{Actores del sistema}
Los actores que interactúan con el sistema Nozama se describen a continuación:

\begin{longtable}{|l|p{7.5cm}|p{5.5cm}|}
    \hline
    \textbf{Actor}                                                                          & \textbf{Descripción} & \textbf{Permisos principales} \\ \hline
    \endfirsthead
    \hline
    \textbf{Actor}                                                                          & \textbf{Descripción} & \textbf{Permisos principales} \\ \hline
    \endhead

    Invitado                                                                                &
    Usuario sin cuenta que puede navegar el catálogo, registrarse e iniciar sesión.         &
    \begin{itemize}[leftmargin=*, nosep]
        \item Navegar catálogo público.
        \item Buscar y filtrar productos.
        \item Registrarse para crear cuenta.
    \end{itemize}                                                                                                        \\ \hline

    Cliente                                                                                 &
    Usuario registrado que hereda las funciones del invitado y posee persistencia de datos. &
    \begin{itemize}[leftmargin=*, nosep]
        \item Gestionar carrito de compras.
        \item Realizar pagos y confirmar órdenes.
        \item Consultar historial de órdenes.
    \end{itemize}                                                                                                   \\ \hline

    Cliente VIP                                                                             &
    Usuario registrado con beneficios adicionales detectados por el sistema.                &
    \begin{itemize}[leftmargin=*, nosep]
        \item Todas las funciones del cliente.
        \item Descuento automático del 10\%.
    \end{itemize}                                                                                                      \\ \hline

    Administrador                                                                           &
    Actor con acceso exclusivo al panel de gestión y control de inventario.                 &
    \begin{itemize}[leftmargin=*, nosep]
        \item Gestión de productos y categorías.
        \item Administración de usuarios.
        \item Gestión de promociones.
    \end{itemize}                                                                                                    \\ \hline

    Sistema                                                                                 &
    Componente interno que ejecuta lógica, validaciones y persiste datos.                   &
    \begin{itemize}[leftmargin=*, nosep]
        \item Actualización de stock y cálculos.
        \item Validación de reglas de negocio.
        \item Cifrado y seguridad de datos.
    \end{itemize}
    \\ \hline
\end{longtable}

% ============================
% REGLAS DE NEGOCIO
% ============================

\subsection{Reglas de negocio}

\begin{enumerate}[label=RN\arabic*:]
    \item El correo electrónico de cada usuario debe ser único en la base de datos.
    \item Las contraseñas deben cifrarse antes de persistirse en la base de datos.
    \item Solo usuarios activos pueden iniciar sesión en el sistema.
    \item La sesión expira tras 30 minutos de inactividad.
    \item Si un usuario inicia sesión en otro dispositivo, la sesión anterior se invalida automáticamente.
    \item Si un usuario inicia sesión desde un dispositivo mientras tiene otra sesión activa, la sesión anterior se invalida automáticamente sin notificación al usuario previo.
    \item La desactivación de un usuario no elimina ni altera sus registros financieros u órdenes previas.
    \item Usuarios desactivados no pueden iniciar sesión ni realizar nuevas transacciones.
    \item La búsqueda de productos se realiza únicamente sobre productos activos.
    \item Cada categoría debe tener un nombre único.
    \item La baja de categorías es lógica (estado inactivo), no física.
    \item Categorías dadas de baja pueden ser reactivadas.
    \item La baja lógica de categorías no elimina ni altera referencias históricas de productos asociados.
    \item Cada producto debe tener un código UPC único.
    \item La baja de productos es lógica (estado inactivo), no física.
    \item Productos dados de baja pueden ser reactivados.
    \item La baja lógica de productos no elimina ni altera órdenes históricas.
    \item Un cliente no puede agregar más unidades al carrito que las disponibles en stock.
    \item El límite máximo de unidades por producto en el carrito es el menor valor entre el stock disponible y 10 unidades.
    \item El carrito debe pertenecer al cliente autenticado.
    \item No se pueden confirmar órdenes con carrito vacío.
    \item El stock debe validarse nuevamente al momento de confirmar una orden.
    \item El stock se reserva y descuenta únicamente al confirmar la orden.
    \item Las promociones deben tener fecha de inicio y fin claramente definidas.
    \item Se aplican promociones vigentes y beneficios VIP automáticamente.
    \item Una promoción activa de categoría tiene prioridad sobre el precio estándar del producto, pero no se acumula con el descuento VIP.
    \item[RN26b:] Si un producto tiene simultáneamente una promoción de producto y una promoción de categoría activas, se aplica únicamente el descuento de mayor porcentaje; no se acumulan entre sí ni con el descuento VIP.
    \item Solo administradores pueden gestionar productos, categorías, usuarios y promociones.
    \item Todas las rutas administrativas deben validarse mediante middleware de autorización.
    \item Seguridad avanzada: contraseñas cifradas con algoritmos seguros (bcrypt/Argon2), limitación de intentos de login, validación y sanitización de entradas contra inyecciones.
\end{itemize}


% ============================
% REQUISITOS FUNCIONALES
% ============================

\subsection{Requisitos funcionales}

\begin{longtable}{|l|p{5.5cm}|p{3cm}|p{1.5cm}|p{4cm}|}
    \hline
    \textbf{ID} & \textbf{Descripción}                                                                       & \textbf{Actor principal} & \textbf{Prior.} & \textbf{Criterio de aceptación} \\ \hline
    \endfirsthead
    \hline
    \textbf{ID} & \textbf{Descripción}                                                                       & \textbf{Actor principal} & \textbf{Prior.} & \textbf{Criterio de aceptación} \\ \hline
    \endhead

    RF01        & Registrar nuevos clientes con nombre, correo electrónico y contraseña.                     & Invitado                 & Alta            &
    El sistema acepta correos con formato RFC 5322, rechaza duplicados con mensaje específico y exige contraseñas de mínimo 8 caracteres con al menos un número.                            \\ \hline

    RF02        & Autenticar usuarios (clientes y administradores) mediante correo electrónico y contraseña. & Cliente / Administrador  & Alta            &
    El sistema bloquea el acceso tras 5 intentos fallidos consecutivos; redirige según rol tras autenticación exitosa.                                                                      \\ \hline

    RF02b       & Cerrar sesión de usuario autenticado.                                                      & Cliente / Administrador  & Alta            &
    El sistema invalida la sesión en BD y redirige al inicio; el botón de retroceso no permite acceder a rutas protegidas.                                                                  \\ \hline

    RF03a       & Mostrar catálogo de productos organizados por categorías a cualquier usuario.              & Invitado / Cliente       & Alta            &
    Se muestran únicamente productos con \texttt{activo = True}; la paginación muestra máximo 20 productos por página.                                                                      \\ \hline

    RF03b       & Permitir búsqueda de productos por término en nombre o descripción.                        & Invitado / Cliente       & Alta            &
    La búsqueda es insensible a mayúsculas; retorna resultados en menos de 2 s; muestra mensaje si no hay coincidencias.                                                                    \\ \hline

    RF03c       & Permitir filtrado de productos por categoría activa.                                       & Invitado / Cliente       & Alta            &
    Solo se muestran categorías con \texttt{activo = True}; al seleccionar una categoría inexistente redirige al catálogo completo.                                                         \\ \hline

    RF04        & Permitir al cliente gestionar el carrito de compras (agregar, actualizar, eliminar).       & Cliente                  & Alta            &
    No se puede agregar más unidades que el stock disponible ni más de 10 por producto; el carrito persiste en sesión.                                                                      \\ \hline

    RF05        & Calcular el total de la orden considerando precio unitario, cantidad y descuentos.         & Cliente / Sistema        & Alta            &
    El total refleja promociones vigentes y descuento VIP (no acumulables); el desglose es visible antes de confirmar.                                                                      \\ \hline

    RF06        & Registrar la orden y descontar stock al confirmar el pago.                                 & Cliente / Sistema        & Alta            &
    El stock se descuenta atómicamente al confirmar; si falla la transacción, el stock no se modifica.                                                                                      \\ \hline

    RF07        & Consultar historial de órdenes de compra anteriores.                                       & Cliente                  & Media           &
    Solo se muestran órdenes del cliente autenticado; se puede filtrar por estado; el detalle incluye productos, cantidades y totales.                                                      \\ \hline

    RF08        & Administrador gestiona productos y categorías (altas, bajas, modificaciones).              & Administrador            & Alta            &
    La baja es lógica (no física); un producto dado de baja no aparece en catálogo ni en búsquedas; puede reactivarse.                                                                      \\ \hline

    RF09        & Administrador gestiona cuentas de usuario (activar/desactivar, cambiar tipo de cliente).   & Administrador            & Alta            &
    Un usuario desactivado no puede iniciar sesión; sus órdenes históricas no se alteran; el cambio de tipo actualiza descuentos en tiempo real.                                            \\ \hline

    RF10        & Soportar clientes estándar y VIP con descuento automático del 10\%.                        & Sistema                  & Alta            &
    El tipo VIP lo asigna manualmente el administrador (RF09); el descuento se aplica automáticamente en RF05 cuando el tipo es VIP.                                                        \\ \hline

    RF11        & Crear promociones de descuento por producto o categoría con fechas de vigencia.            & Administrador            & Media           &
    La fecha de inicio debe ser anterior a la de fin; una promoción fuera de vigencia no se aplica aunque esté activa; no se acumula con descuento VIP.                                     \\ \hline

    RF12        & Captura de productos en formato UPC (código de barras, nombre, descripción, cantidad).     & Administrador            & Alta            &
    El sistema rechaza UPCs duplicados con mensaje específico; el campo UPC acepta exactamente 12 dígitos numéricos.                                                                        \\ \hline
\end{longtable}

% ============================
% REQUISITOS NO FUNCIONALES
% ============================

\subsection{Requisitos no funcionales}

\begin{longtable}{|l|p{7cm}|p{3cm}|p{3.5cm}|}
    \hline
    \textbf{ID} & \textbf{Descripción}                                                                                                  & \textbf{Categoría}     & \textbf{Métrica / criterio} \\ \hline
    \endfirsthead
    \hline
    \textbf{ID} & \textbf{Descripción}                                                                                                  & \textbf{Categoría}     & \textbf{Métrica / criterio} \\ \hline
    \endhead

    RNF01       & Tiempo de respuesta menor a 2 segundos con hasta 100 usuarios concurrentes.                                           & Rendimiento            &
    Medido con herramienta de carga (e.g., Locust) sobre rutas críticas: catálogo, carrito y checkout.                                                                                         \\ \hline

    RNF02       & Contraseñas cifradas con hash seguro (bcrypt o Argon2). Rutas de administración protegidas por autenticación y roles. & Seguridad              &
    Acceso a rutas \texttt{/admin/*} sin sesión retorna HTTP 403; ninguna contraseña se almacena en texto plano.                                                                               \\ \hline

    RNF03       & Interfaz responsiva mediante Bootstrap 5, funcional en dispositivos móviles y escritorio.                             & Usabilidad             &
    Compatible con Chrome 120+, Firefox 120+ y Safari 17+; resoluciones mínimas: 375px (móvil) y 1280px (escritorio).                                                                          \\ \hline

    RNF04       & Disponibilidad mínima del 95\% durante el período de evaluación operativa.                                            & Disponibilidad         &
    Calculado como: (tiempo total $-$ tiempo de caída) / tiempo total $\geq$ 0.95 en el período de evaluación.                                                                                 \\ \hline

    RNF05       & Código estructurado bajo patrón MVC (Blueprints, SQLAlchemy, Jinja2) para facilitar mantenibilidad.                   & Mantenibilidad         &
    Cada módulo funcional en un Blueprint independiente; sin lógica de negocio en las plantillas Jinja2.                                                                                       \\ \hline

    RNF06       & Compatible con entornos que dispongan de Python 3.x, Flask y MySQL.                                                   & Portabilidad           &
    El proyecto incluye \texttt{requirements.txt} y script de inicialización de BD; ejecutable en Ubuntu 22.04+.                                                                               \\ \hline

    RNF07       & Registro de errores en logs con timestamp y despliegue de mensajes genéricos al usuario.                              & Confiabilidad          &
    Formato de log: \texttt{[YYYY-MM-DD HH:MM:SS] [LEVEL] mensaje}; niveles mínimos registrados: WARNING y ERROR.                                                                              \\ \hline

    RNF08       & Comunicación visual consistente: notificaciones uniformes para confirmar acciones del usuario.                        & Experiencia de usuario &
    Todas las acciones CRUD muestran mensaje flash de confirmación o error; el tiempo de visualización es de al menos 3 segundos.                                                              \\ \hline

    RNF09       & El sistema opera en idioma español (es-MX) como único idioma soportado.                                               & Localización           &
    Todos los mensajes de interfaz, errores y notificaciones están en español; no se requiere soporte multilingüe.                                                                             \\ \hline
\end{longtable}

\section{Casos de uso}

% ============================
% CU01
% ============================

\subsection{CU01: Registrar cuenta}
\subsubsection*{Objetivo}
Permitir a un invitado crear una cuenta en el sistema, validando datos y generando sesión automática.

\subsubsection*{Actores}
\begin{longtable}{|l|p{10cm}|}
    \hline
    Invitado & Proporciona datos personales para crear una cuenta.      \\ \hline
    Sistema  & Valida datos, cifra contraseña y persiste usuario en BD. \\ \hline
\end{longtable}

\subsubsection*{Condiciones iniciales}
\begin{itemize}[label=--]
    \item El usuario no tiene cuenta registrada.
    \item No hay sesión activa.
    \item BD y servicio web disponibles.
\end{itemize}

\subsubsection*{Condiciones de salida}
\begin{longtable}{|l|p{10cm}|}
    \hline
    Éxito             & Usuario creado, sesión iniciada, redirección al catálogo. \\ \hline
    Correo duplicado  & Error en campo, sugerencia de iniciar sesión.             \\ \hline
    Datos inválidos   & Errores específicos por campo.                            \\ \hline
    Error del sistema & Mensaje genérico, registro en log.                        \\ \hline
\end{longtable}

\subsubsection*{Flujo normal}
\begin{enumerate}[label=FN-\arabic*]
    \item Invitado accede a \texttt{/register}.
    \item Completa formulario con nombre, correo y contraseña.
    \item Sistema valida campos.
    \item Sistema verifica correo no duplicado.
    \item Sistema cifra contraseña y guarda usuario.
    \item Sistema inicia sesión y redirige a \texttt{/catalog}.
\end{enumerate}

\subsubsection*{Flujos alternativos}
\begin{enumerate}[label=FA-\arabic*]
    \item Correo ya registrado: sistema muestra error y sugiere iniciar sesión.
    \item Datos inválidos: sistema muestra errores específicos y mantiene formulario visible.
\end{enumerate}

\subsubsection*{Flujos excepcionales}
\begin{enumerate}[label=FE-\arabic*]
    \item Error en BD: sistema captura excepción, registra log y muestra mensaje genérico.
    \item Timeout: invitado no recibe respuesta, sistema evita duplicados.
\end{enumerate}

% ============================
% CU02
% ============================

\subsection{CU02: Iniciar sesión}
\subsubsection*{Objetivo}
Permitir a clientes y administradores autenticarse en el sistema.

\subsubsection*{Actores}
\begin{longtable}{|l|p{10cm}|}
    \hline
    Cliente registrado & Ingresa credenciales para acceder a su cuenta.    \\ \hline
    Administrador      & Se autentica y accede al panel de administración. \\ \hline
    Sistema            & Verifica credenciales y activa sesión según rol.  \\ \hline
\end{longtable}

\subsubsection*{Condiciones iniciales}
\begin{itemize}[label=--]
    \item Usuario sin sesión activa.
    \item Cuenta registrada en el sistema.
    \item BD y servicio web disponibles.
\end{itemize}

\subsubsection*{Condiciones de salida}
\begin{longtable}{|l|p{10cm}|}
    \hline
    Éxito (cliente)          & Sesión iniciada, redirección al catálogo.                \\ \hline
    Éxito (administrador)    & Sesión iniciada, redirección al panel de administración. \\ \hline
    Credenciales incorrectas & Error genérico sin revelar campo fallido.                \\ \hline
    Cuenta desactivada       & Mensaje informativo y sugerencia de soporte.             \\ \hline
    Error del sistema        & Mensaje genérico y registro en log.                      \\ \hline
\end{longtable}

\subsubsection*{Flujo normal}
\begin{enumerate}[label=FN-\arabic*]
    \item Usuario accede a \texttt{/login}.
    \item Ingresa correo y contraseña.
    \item Sistema busca usuario en BD.
    \item Sistema verifica contraseña.
    \item Sistema comprueba estado activo.
    \item Sistema inicia sesión y redirige según rol.
\end{enumerate}

\subsubsection*{Flujos alternativos}
\begin{enumerate}[label=FA-\arabic*]
    \item Credenciales incorrectas: sistema muestra error genérico.
    \item Cuenta desactivada: sistema informa y sugiere soporte.
\end{enumerate}

\subsubsection*{Flujos excepcionales}
\begin{enumerate}[label=FE-\arabic*]
    \item Error en BD: sistema captura excepción y muestra mensaje genérico.
    \item Sesión activa: sistema detecta sesión y redirige sin mostrar formulario.
\end{enumerate}

% ============================
% CU02b
% ============================

\subsection{CU02b: Cerrar sesión}
\subsubsection*{Objetivo}
Permitir a clientes y administradores cerrar su sesión activa de forma segura.

\subsubsection*{Actores}
\begin{longtable}{|l|p{10cm}|}
    \hline
    Cliente / Administrador & Solicita cerrar su sesión activa.              \\ \hline
    Sistema                 & Invalida la sesión en BD y redirige al inicio. \\ \hline
\end{longtable}

\subsubsection*{Condiciones iniciales}
\begin{itemize}[label=--]
    \item El usuario tiene sesión activa en el sistema.
    \item BD y servicio web disponibles.
\end{itemize}

\subsubsection*{Condiciones de salida}
\begin{longtable}{|l|p{10cm}|}
    \hline
    Éxito             & Sesión invalidada; redirección a \texttt{/login}.                         \\ \hline
    Error del sistema & Sesión no invalidada; mensaje genérico y registro en log.                 \\ \hline
\end{longtable}

\subsubsection*{Flujo normal}
\begin{enumerate}[label=FN-\arabic*]
    \item Usuario hace clic en "Cerrar sesión".
    \item Sistema invalida el token/sesión en la BD.
    \item Sistema redirige a \texttt{/login}.
    \item Intentar navegar hacia atrás retorna HTTP 403 o redirige a \texttt{/login}.
\end{enumerate}

\subsubsection*{Flujos excepcionales}
\begin{enumerate}[label=FE-\arabic*]
    \item Error en BD: sistema captura excepción, registra log; la sesión se invalida localmente de todos modos.
\end{enumerate}

% ============================
% CU03
% ============================

\subsection{CU03: Ver catálogo}
\subsubsection*{Objetivo}
Permitir al cliente navegar el catálogo de productos activos, aplicar filtros por categoría y consultar el detalle de productos.

\subsubsection*{Actores}
\begin{longtable}{|l|p{10cm}|}
    \hline
    Cliente & Navega el catálogo, filtra por categoría y consulta detalle de productos.                \\ \hline
    Sistema & Recupera y muestra los productos activos desde la base de datos según filtros aplicados. \\ \hline
\end{longtable}

\subsubsection*{Condiciones iniciales}
\begin{itemize}[label=--]
    \item El cliente puede estar autenticado o no (el catálogo es público).
    \item Existen productos registrados y marcados como activos en la base de datos.
    \item BD y servicio web disponibles.
\end{itemize}

\subsubsection*{Condiciones de salida}
\begin{longtable}{|l|p{10cm}|}
    \hline
    Éxito             & Se muestran productos activos con imagen, nombre, precio y stock.             \\ \hline
    Sin resultados    & No hay productos que coincidan con el filtro; se muestra mensaje informativo. \\ \hline
    Catálogo vacío    & No existen productos activos en la BD; se muestra estado vacío con mensaje.   \\ \hline
    Error del sistema & Mensaje genérico y registro en log.                                           \\ \hline
\end{longtable}

\subsubsection*{Flujo normal}
\begin{enumerate}[label=FN-\arabic*]
    \item Cliente accede a \texttt{/catalog}.
    \item Sistema consulta todos los productos con \texttt{activo = True} y los devuelve paginados.
    \item Sistema renderiza tarjetas de producto con imagen, nombre, precio, categoría y botón “Agregar al carrito”.
    \item Cliente selecciona una categoría del menú de filtros.
    \item Sistema filtra productos por categoría y actualiza la vista.
    \item Cliente hace clic en un producto para ver su detalle.
    \item Sistema renderiza la página de detalle en \texttt{/product/<id>}.
\end{enumerate}

\subsubsection*{Flujos alternativos}
\begin{enumerate}[label=FA-\arabic*]
    \item Búsqueda sin resultados: sistema muestra mensaje “No se encontraron productos en esta categoría” con opción de limpiar filtros.
    \item Producto sin stock: sistema muestra etiqueta “Agotado” y deshabilita botón “Agregar al carrito”.
\end{enumerate}

\subsubsection*{Flujos excepcionales}
\begin{enumerate}[label=FE-\arabic*]
    \item Error en BD: sistema captura excepción, registra log y muestra mensaje genérico.
    \item Producto inexistente o inactivo: sistema retorna error 404 con mensaje “Este producto no está disponible” y enlace de regreso al catálogo.
\end{enumerate}

% ============================
% CU04
% ============================

\subsection{CU04: Buscar productos}
\subsection*{Objetivo}
Permitir al usuario localizar productos específicos mediante búsqueda por término dentro del catálogo.

\subsubsection*{Actores}
\begin{longtable}{|l|p{10cm}|}
    \hline
    Invitado/Cliente & Ingresa un término de búsqueda para localizar productos específicos.                       \\ \hline
    Sistema          & Ejecuta la consulta sobre la base de datos y devuelve los productos activos que coincidan. \\ \hline
\end{longtable}

\subsubsection*{Condiciones iniciales}
\begin{itemize}[label=--]
    \item El usuario (invitado o cliente) se encuentra en el catálogo (\texttt{/catalog}).
    \item BD y servicio web disponibles.
    \item El producto buscado debe estar marcado como activo en el sistema.
\end{itemize}

\subsubsection*{Condiciones de salida}
\begin{longtable}{|l|p{10cm}|}
    \hline
    Con resultados    & Se muestran productos activos cuyo nombre o descripción coincidan con el término.    \\ \hline
    Sin resultados    & No hay coincidencias; se muestra mensaje informativo con opción de limpiar búsqueda. \\ \hline
    Error del sistema & No se ejecuta la búsqueda; mensaje genérico y registro en log.                       \\ \hline
\end{longtable}

\subsubsection*{Flujo normal}
\begin{enumerate}[label=FN-\arabic*]
    \item Usuario escribe un término en la barra de búsqueda y presiona Enter o clic en el icono de buscar.
    \item Sistema recibe el parámetro \texttt{q} vía GET en \texttt{/catalog?q=termino}.
    \item Sistema consulta productos activos cuyo nombre o descripción contengan el término (búsqueda insensible a mayúsculas).
    \item Sistema renderiza resultados paginados con el término resaltado en los títulos.
\end{enumerate}

\subsubsection*{Flujos alternativos}
\begin{enumerate}[label=FA-\arabic*]
    \item Sin resultados: la consulta retorna lista vacía; sistema muestra mensaje “No se encontraron productos para \{término\}” con botón “Ver todo el catálogo”.
    \item Búsqueda vacía: sistema detecta que el parámetro \texttt{q} está vacío o contiene solo espacios; redirige al catálogo completo sin aplicar filtro.
\end{enumerate}

\subsubsection*{Flujos excepcionales}
\begin{enumerate}[label=FE-\arabic*]
    \item Error en BD: sistema captura excepción, registra log y muestra mensaje genérico invitando a intentarlo de nuevo.
\end{enumerate}

% ============================
% CU05
% ============================

\subsection{CU05: Filtrar por categoría}
\subsection*{Objetivo}
Permitir al usuario filtrar productos activos por una categoría seleccionada dentro del catálogo.

\subsubsection*{Actores}
\begin{longtable}{|l|p{10cm}|}
    \hline
    Invitado/Cliente & Selecciona una categoría para acotar los productos mostrados en el catálogo.     \\ \hline
    Sistema          & Filtra los productos activos por la categoría seleccionada y actualiza la vista. \\ \hline
\end{longtable}

\subsubsection*{Condiciones iniciales}
\begin{itemize}[label=--]
    \item El usuario (invitado o cliente) se encuentra en el catálogo (\texttt{/catalog}).
    \item Existen categorías activas con al menos un producto activo asociado.
    \item BD y servicio web disponibles.
\end{itemize}

\subsubsection*{Condiciones de salida}
\begin{longtable}{|l|p{10cm}|}
    \hline
    Con resultados        & Se muestran únicamente los productos que pertenecen a la categoría seleccionada. \\ \hline
    Sin resultados        & La categoría no tiene productos activos; se muestra mensaje informativo.         \\ \hline
    Categoría inexistente & El ID de categoría en la URL no existe; se redirige al catálogo completo.        \\ \hline
    Error del sistema     & No se aplica el filtro; mensaje genérico y registro en log.                      \\ \hline
\end{longtable}

\subsubsection*{Flujo normal}
\begin{enumerate}[label=FN-\arabic*]
    \item Usuario hace clic en una categoría del menú de navegación o filtros laterales.
    \item Sistema recibe el parámetro \texttt{categoria\_id} vía GET en \texttt{/catalog?categoria=id}.
    \item Sistema verifica que la categoría exista y esté activa en la BD.
    \item Sistema consulta todos los productos activos con \texttt{categoria\_id} igual al seleccionado.
    \item Sistema renderiza resultados paginados con la categoría activa resaltada en el menú.
\end{enumerate}

\subsubsection*{Flujos alternativos}
\begin{enumerate}[label=FA-\arabic*]
    \item Categoría sin productos activos: la consulta retorna lista vacía; sistema muestra mensaje “Esta categoría no tiene productos disponibles por el momento” con enlace al catálogo completo.
    \item Categoría inexistente en la URL: sistema no encuentra la categoría en la BD; redirige al catálogo completo sin filtro aplicado.
\end{enumerate}

\subsubsection*{Flujos excepcionales}
\begin{enumerate}[label=FE-\arabic*]
    \item Error en BD: sistema captura excepción, registra log y muestra mensaje genérico invitando a intentarlo de nuevo.
\end{enumerate}

% ============================
% CU06
% ============================

\subsection{CU06: Ver detalle de producto}
\subsubsection*{Objetivo}
Mostrar información detallada de un producto seleccionado, incluyendo descripción, precio, stock y promociones vigentes.

\subsubsection*{Actores}
\begin{longtable}{|l|p{10cm}|}
    \hline
    Cliente & Selecciona un producto para consultar su detalle.                                             \\ \hline
    Sistema & Recupera la información completa del producto desde la base de datos y la muestra al cliente. \\ \hline
\end{longtable}

\subsubsection*{Condiciones iniciales}
\begin{itemize}[label=--]
    \item El cliente se encuentra en el catálogo (\texttt{/catalog}).
    \item El producto seleccionado está registrado y marcado como activo en la BD.
    \item BD y servicio web disponibles.
\end{itemize}

\subsubsection*{Condiciones de salida}
\begin{longtable}{|l|p{10cm}|}
    \hline
    Éxito                         & Se muestra la información completa del producto (nombre, descripción, precio, stock, categoría, promociones). \\ \hline
    Producto agotado              & Se muestra detalle con etiqueta “Agotado” y botón de compra deshabilitado.                                    \\ \hline
    Producto inexistente/inactivo & Se retorna error 404 con mensaje informativo.                                                                 \\ \hline
    Error del sistema             & Mensaje genérico y registro en log.                                                                           \\ \hline
\end{longtable}

\subsubsection*{Flujo normal}
\begin{enumerate}[label=FN-\arabic*]
    \item Cliente hace clic en un producto desde el catálogo.
    \item Sistema recibe el \texttt{product\_id} vía GET en \texttt{/product/<id>}.
    \item Sistema consulta la BD y obtiene información completa del producto.
    \item Sistema renderiza la página de detalle con nombre, descripción, precio, stock, categoría y promociones vigentes.
\end{enumerate}

\subsubsection*{Flujos alternativos}
\begin{enumerate}[label=FA-\arabic*]
    \item Producto agotado: sistema detecta stock = 0; muestra etiqueta “Agotado” y deshabilita botón “Agregar al carrito”.
    \item Producto con promoción vigente: sistema detecta promoción activa; muestra precio original tachado y precio con descuento.
\end{enumerate}

\subsubsection*{Flujos excepcionales}
\begin{enumerate}[label=FE-\arabic*]
    \item Producto inexistente o inactivo: sistema no encuentra el producto o \texttt{activo = False}; retorna error 404 con mensaje “Este producto no está disponible” y enlace de regreso al catálogo.
    \item Error en BD: sistema captura excepción, registra log y muestra mensaje genérico invitando a intentarlo de nuevo.
\end{enumerate}

% ============================
% CU07
% ============================

\subsection{CU07: Agregar producto al carrito}
\subsubsection*{Objetivo}
Permitir al cliente agregar productos al carrito de compras, validando stock y reglas de negocio.

\subsubsection*{Actores}
\begin{longtable}{|l|p{10cm}|}
    \hline
    Cliente & Selecciona un producto y lo agrega al carrito.                         \\ \hline
    Sistema & Valida stock, cantidad y sesión activa; actualiza el carrito en la BD. \\ \hline
\end{longtable}

\subsubsection*{Condiciones iniciales}
\begin{itemize}[label=--]
    \item El cliente tiene sesión activa en el sistema.
    \item El producto seleccionado está registrado, activo y con stock disponible.
    \item BD y servicio web disponibles.
\end{itemize}

\subsubsection*{Condiciones de salida}
\begin{longtable}{|l|p{10cm}|}
    \hline
    Éxito                     & Producto agregado al carrito; se muestra confirmación.     \\ \hline
    Producto ya en el carrito & Sistema incrementa cantidad o muestra mensaje informativo. \\ \hline
    Cantidad máxima alcanzada & Sistema muestra error y no permite agregar más unidades.   \\ \hline
    Producto agotado/inactivo & Sistema muestra mensaje de no disponibilidad.              \\ \hline
    Error del sistema         & Mensaje genérico y registro en log.                        \\ \hline
\end{longtable}

\subsubsection*{Flujo normal}
\begin{enumerate}[label=FN-\arabic*]
    \item Cliente hace clic en “Agregar al carrito” desde el catálogo o detalle de producto.
    \item Sistema verifica que el cliente tenga sesión activa.
    \item Sistema valida que el producto esté activo y con stock disponible.
    \item Sistema agrega el producto al carrito en la BD con cantidad inicial = 1.
    \item Sistema muestra confirmación y actualiza el ícono del carrito.
\end{enumerate}

\subsubsection*{Flujos alternativos}
\begin{enumerate}[label=FA-\arabic*]
    \item Producto ya en el carrito: sistema detecta existencia previa; incrementa cantidad o muestra mensaje “Este producto ya está en tu carrito”.
    \item Cantidad máxima alcanzada: sistema detecta límite de unidades por producto; muestra error y no permite agregar más.
\end{enumerate}

\subsubsection*{Flujos excepcionales}
\begin{enumerate}[label=FE-\arabic*]
    \item Sin sesión activa: sistema detecta que el cliente no está autenticado; redirige a \texttt{/login}.
    \item Producto agotado o inactivo: sistema detecta stock = 0 o \texttt{activo = False}; muestra mensaje de no disponibilidad.
    \item Error en BD: sistema captura excepción, registra log y muestra mensaje genérico invitando a intentarlo de nuevo.
\end{enumerate}

% ============================
% CU08
% ============================

\subsection{CU08: Gestionar carrito}
\subsubsection*{Objetivo}
Permitir al cliente modificar las cantidades de los productos en su carrito, eliminar ítems o vaciarlo por completo, validando stock y reglas de negocio.

\subsubsection*{Actores}
\begin{longtable}{|l|p{10cm}|}
    \hline
    Cliente & Gestiona los productos de su carrito (actualizar cantidad, eliminar ítems, vaciar carrito). \\ \hline
    Sistema & Valida stock, cantidades y sesión activa; actualiza el carrito en la BD.                    \\ \hline
\end{longtable}

\subsubsection*{Condiciones iniciales}
\begin{itemize}[label=--]
    \item El cliente tiene sesión activa en el sistema.
    \item El carrito contiene al menos un producto.
    \item BD y servicio web disponibles.
\end{itemize}

\subsubsection*{Condiciones de salida}
\begin{longtable}{|l|p{10cm}|}
    \hline
    Éxito              & Carrito actualizado correctamente (cantidad modificada, ítem eliminado o carrito vaciado). \\ \hline
    Cantidad inválida  & Sistema muestra error y no aplica cambios.                                                 \\ \hline
    Stock insuficiente & Sistema muestra mensaje informativo y ajusta cantidad máxima posible.                      \\ \hline
    Carrito vacío      & Tras eliminar todos los ítems, se muestra estado vacío.                                    \\ \hline
    Error del sistema  & Mensaje genérico y registro en log.                                                        \\ \hline
\end{longtable}

\subsubsection*{Flujo normal}
\begin{enumerate}[label=FN-\arabic*]
    \item Cliente accede a su carrito desde la barra de navegación.
    \item Cliente selecciona acción: actualizar cantidad, eliminar ítem o vaciar carrito.
    \item Sistema valida sesión activa y existencia de ítems en el carrito.
    \item Sistema aplica la acción solicitada en la BD.
    \item Sistema muestra confirmación y actualiza la vista del carrito.
\end{enumerate}

\subsubsection*{Flujos alternativos}
\begin{enumerate}[label=FA-\arabic*]
    \item Cantidad inválida: sistema detecta número negativo o no numérico; muestra error y mantiene cantidad anterior.
    \item Stock insuficiente: sistema detecta que la cantidad solicitada excede stock disponible; ajusta cantidad máxima y muestra mensaje informativo.
    \item Carrito vacío tras eliminar: sistema detecta que no quedan ítems; muestra estado vacío con mensaje “Tu carrito está vacío”.
    \item Vaciar carrito completo: cliente selecciona opción “Vaciar carrito”; sistema elimina todos los ítems y muestra estado vacío.
\end{enumerate}

\subsubsection*{Flujos excepcionales}
\begin{enumerate}[label=FE-\arabic*]
    \item Sin sesión activa: sistema detecta cliente no autenticado; redirige a \texttt{/login}.
    \item Ítem no pertenece al cliente: sistema detecta inconsistencia; muestra error y no aplica cambios.
    \item Error en BD: sistema captura excepción, registra log y muestra mensaje genérico invitando a intentarlo de nuevo.
\end{enumerate}

% ============================
% CU09
% ============================

\subsection{CU09: Realizar pago y confirmar orden}
\subsubsection*{Objetivo}
Permitir al cliente confirmar su orden y realizar el pago, validando stock, sesión activa y reglas de negocio.

\subsubsection*{Actores}
\begin{longtable}{|l|p{10cm}|}
    \hline
    Cliente & Selecciona método de pago y confirma la orden.                            \\ \hline
    Sistema & Valida stock, sesión activa y método de pago; registra la orden en la BD. \\ \hline
\end{longtable}

\subsubsection*{Condiciones iniciales}
\begin{itemize}[label=--]
    \item El cliente tiene sesión activa en el sistema.
    \item El carrito contiene al menos un producto con stock disponible.
    \item BD y servicio web disponibles.
\end{itemize}

\subsubsection*{Condiciones de salida}
\begin{longtable}{|l|p{10cm}|}
    \hline
    Éxito                          & Orden confirmada y registrada; se muestra comprobante de pago. \\ \hline
    Cliente VIP                    & Se aplican beneficios especiales (ej. descuentos adicionales). \\ \hline
    Promoción vigente              & Se aplican descuentos sobre productos con promoción activa.    \\ \hline
    Método de pago no seleccionado & Sistema muestra error y solicita elegir método de pago.        \\ \hline
    Stock insuficiente             & Sistema muestra error y no confirma orden.                     \\ \hline
    Carrito vacío                  & Sistema muestra mensaje informativo y no permite confirmar.    \\ \hline
    Sin sesión activa              & Sistema redirige a \texttt{/login}.                            \\ \hline
    Error del sistema              & Mensaje genérico y registro en log.                            \\ \hline
\end{longtable}

\subsubsection*{Flujo normal}
\begin{enumerate}[label=FN-\arabic*]
    \item Cliente accede a su carrito y selecciona “Confirmar orden”.
    \item Sistema valida sesión activa y existencia de productos en el carrito.
    \item Cliente selecciona método de pago simulado (p.~ej., ``Pago en efectivo'' o ``Transferencia bancaria''); no se integran pasarelas de pago externas en esta versión.
    \item Sistema valida stock de cada producto antes de confirmar.
    \item Sistema aplica promociones vigentes y beneficios VIP si corresponden.
    \item Sistema registra la orden en la BD y procesa el pago.
    \item Sistema muestra comprobante de pago y redirige a página de confirmación.
\end{enumerate}

\subsubsection*{Flujos alternativos}
\begin{enumerate}[label=FA-\arabic*]
    \item Cliente VIP: sistema detecta atributo VIP y aplica beneficios adicionales.
    \item Promoción vigente: sistema aplica descuentos sobre productos con promoción activa.
    \item Método de pago no seleccionado: sistema muestra error y solicita elegir método antes de continuar.
\end{enumerate}

\subsubsection*{Flujos excepcionales}
\begin{enumerate}[label=FE-\arabic*]
    \item Stock insuficiente: sistema detecta falta de stock al confirmar; muestra error y no registra orden.
    \item Carrito vacío: sistema detecta que no hay productos; muestra mensaje informativo y no permite confirmar.
    \item Sin sesión activa: sistema detecta cliente no autenticado; redirige a \texttt{/login}.
    \item Error en BD: sistema captura excepción, registra log y muestra mensaje genérico invitando a intentarlo de nuevo.
\end{enumerate}

% ============================
% CU10
% ============================

\subsection{CU10: Consultar historial de órdenes}
\subsubsection*{Objetivo}
Permitir al cliente consultar su historial de órdenes, ver detalles de cada una y aplicar filtros por estado.

\subsubsection*{Actores}
\begin{longtable}{|l|p{10cm}|}
    \hline
    Cliente & Consulta el historial de órdenes y accede al detalle de cada una.   \\ \hline
    Sistema & Recupera las órdenes registradas en la BD y las muestra al cliente. \\ \hline
\end{longtable}

\subsubsection*{Condiciones iniciales}
\begin{itemize}[label=--]
    \item El cliente tiene sesión activa en el sistema.
    \item Existen órdenes registradas en la BD asociadas al cliente.
    \item BD y servicio web disponibles.
\end{itemize}

\subsubsection*{Condiciones de salida}
\begin{longtable}{|l|p{10cm}|}
    \hline
    Éxito                         & Se muestra listado de órdenes con opción de ver detalle. \\ \hline
    Sin órdenes registradas       & Sistema muestra mensaje informativo.                     \\ \hline
    Filtro aplicado               & Sistema muestra órdenes según estado seleccionado.       \\ \hline
    Orden no pertenece al cliente & Sistema muestra error y no permite acceso.               \\ \hline
    Sin sesión activa             & Sistema redirige a \texttt{/login}.                      \\ \hline
    Error del sistema             & Mensaje genérico y registro en log.                      \\ \hline
\end{longtable}

\subsubsection*{Flujo normal}
\begin{enumerate}[label=FN-\arabic*]
    \item Cliente accede a la sección “Historial de órdenes”.
    \item Sistema valida sesión activa y consulta órdenes asociadas al cliente.
    \item Sistema muestra listado de órdenes con fecha, estado y total.
    \item Cliente selecciona una orden para ver detalle.
    \item Sistema muestra detalle de la orden: productos, cantidades, precios y estado.
\end{enumerate}

\subsubsection*{Flujos alternativos}
\begin{enumerate}[label=FA-\arabic*]
    \item Sin órdenes registradas: sistema detecta lista vacía; muestra mensaje “No tienes órdenes registradas”.
    \item Filtrar por estado: cliente selecciona estado (ej. pendiente, entregada); sistema muestra solo órdenes con ese estado.
\end{enumerate}

\subsubsection*{Flujos excepcionales}
\begin{enumerate}[label=FE-\arabic*]
    \item Orden no pertenece al cliente: sistema detecta intento de acceso a orden ajena; muestra error y no permite visualizar detalle.
    \item Sin sesión activa: sistema detecta cliente no autenticado; redirige a \texttt{/login}.
    \item Error en BD: sistema captura excepción, registra log y muestra mensaje genérico invitando a intentarlo de nuevo.
\end{enumerate}

% ============================
% CU11
% ============================

\subsection{CU11: Gestionar productos}
\subsubsection*{Objetivo}
Permitir al administrador dar de alta, modificar, dar de baja o reactivar productos en el sistema, validando datos y reglas de negocio.

\subsubsection*{Actores}
\begin{longtable}{|l|p{10cm}|}
    \hline
    Administrador & Gestiona los productos del catálogo (alta, modificación, baja lógica, reactivación). \\ \hline
    Sistema       & Valida datos, persiste cambios en la BD y aplica reglas de negocio.                  \\ \hline
\end{longtable}

\subsubsection*{Condiciones iniciales}
\begin{itemize}[label=--]
    \item El administrador tiene sesión activa en el sistema.
    \item BD y servicio web disponibles.
\end{itemize}

\subsubsection*{Condiciones de salida}
\begin{longtable}{|l|p{10cm}|}
    \hline
    Éxito                              & Producto creado, modificado, dado de baja o reactivado correctamente. \\ \hline
    UPC duplicado                      & Sistema muestra error y no permite alta.                              \\ \hline
    Datos inválidos                    & Sistema muestra errores específicos por campo.                        \\ \hline
    Sin sesión o sin rol administrador & Sistema redirige a \texttt{/login} o muestra error de permisos.       \\ \hline
    Error del sistema                  & Mensaje genérico y registro en log.                                   \\ \hline
\end{longtable}

\subsubsection*{Flujo normal}
\begin{enumerate}[label=FN-\arabic*]
    \item Administrador accede a la sección “Gestión de productos”.
    \item Selecciona acción: alta, modificación, baja lógica o reactivación.
    \item Sistema valida sesión activa y rol administrador.
    \item Administrador ingresa o modifica datos del producto (nombre, descripción, precio, stock, categoría, UPC).
    \item Sistema valida datos y persiste cambios en la BD.
    \item Sistema muestra confirmación y actualiza el catálogo.
\end{enumerate}

\subsubsection*{Flujos alternativos}
\begin{enumerate}[label=FA-\arabic*]
    \item UPC duplicado: sistema detecta código repetido; muestra error y no permite alta.
    \item Datos inválidos: sistema detecta campos vacíos o incorrectos; muestra errores específicos y mantiene formulario visible.
    \item Reactivar producto dado de baja: administrador selecciona producto inactivo; sistema cambia estado a activo y lo muestra en catálogo.
\end{enumerate}

\subsubsection*{Flujos excepcionales}
\begin{enumerate}[label=FE-\arabic*]
    \item Sin sesión o sin rol administrador: sistema detecta falta de permisos; redirige a \texttt{/login} o muestra error de autorización.
    \item Error en BD: sistema captura excepción, registra log y muestra mensaje genérico invitando a intentarlo de nuevo.
\end{enumerate}

% ============================
% CU12
% ============================

\subsection{CU12: Gestionar categorías}
\subsubsection*{Objetivo}
Permitir al administrador dar de alta, modificar, dar de baja o reactivar categorías en el sistema, validando datos y reglas de negocio.

\subsubsection*{Actores}
\begin{longtable}{|l|p{10cm}|}
    \hline
    Administrador & Gestiona las categorías del catálogo (alta, modificación, baja lógica, reactivación). \\ \hline
    Sistema       & Valida datos, persiste cambios en la BD y aplica reglas de negocio.                   \\ \hline
\end{longtable}

\subsubsection*{Condiciones iniciales}
\begin{itemize}[label=--]
    \item El administrador tiene sesión activa en el sistema.
    \item BD y servicio web disponibles.
\end{itemize}

\subsubsection*{Condiciones de salida}
\begin{longtable}{|l|p{10cm}|}
    \hline
    Éxito                              & Categoría creada, modificada, dada de baja o reactivada correctamente. \\ \hline
    Nombre duplicado                   & Sistema muestra error y no permite alta.                               \\ \hline
    Categoría con productos activos    & Sistema muestra error y no permite baja lógica.                        \\ \hline
    Sin sesión o sin rol administrador & Sistema redirige a \texttt{/login} o muestra error de permisos.        \\ \hline
    Error del sistema                  & Mensaje genérico y registro en log.                                    \\ \hline
\end{longtable}

\subsubsection*{Flujo normal}
\begin{enumerate}[label=FN-\arabic*]
    \item Administrador accede a la sección “Gestión de categorías”.
    \item Selecciona acción: alta, modificación, baja lógica o reactivación.
    \item Sistema valida sesión activa y rol administrador.
    \item Administrador ingresa o modifica datos de la categoría (nombre, descripción).
    \item Sistema valida datos y persiste cambios en la BD.
    \item Sistema muestra confirmación y actualiza el catálogo.
\end{enumerate}

\subsubsection*{Flujos alternativos}
\begin{enumerate}[label=FA-\arabic*]
    \item Nombre duplicado: sistema detecta nombre repetido; muestra error y no permite alta.
    \item Categoría con productos activos: sistema detecta que la categoría tiene productos activos; muestra error y no permite baja lógica.
    \item Reactivar categoría: administrador selecciona categoría inactiva; sistema cambia estado a activo y la muestra en catálogo.
\end{enumerate}

\subsubsection*{Flujos excepcionales}
\begin{enumerate}[label=FE-\arabic*]
    \item Sin sesión o sin rol administrador: sistema detecta falta de permisos; redirige a \texttt{/login} o muestra error de autorización.
    \item Error en BD: sistema captura excepción, registra log y muestra mensaje genérico invitando a intentarlo de nuevo.
\end{enumerate}

% ============================
% CU13
% ============================

\subsection{CU13: Gestionar usuarios}
\subsubsection*{Objetivo}
Permitir al administrador consultar, modificar roles, desactivar o reactivar usuarios en el sistema, validando permisos y reglas de negocio.

\subsubsection*{Actores}
\begin{longtable}{|l|p{10cm}|}
    \hline
    Administrador & Gestiona los usuarios del sistema (consulta, modificación de roles, desactivación, reactivación). \\ \hline
    Sistema       & Valida datos, persiste cambios en la BD y aplica reglas de negocio.                               \\ \hline
\end{longtable}

\subsubsection*{Condiciones iniciales}
\begin{itemize}[label=--]
    \item El administrador tiene sesión activa en el sistema.
    \item BD y servicio web disponibles.
    \item Existen usuarios registrados en la BD.
\end{itemize}

\subsubsection*{Condiciones de salida}
\begin{longtable}{|l|p{10cm}|}
    \hline
    Éxito                              & Usuario consultado, rol modificado, desactivado o reactivado correctamente. \\ \hline
    Usuario inexistente                & Sistema muestra error y no permite acción.                                  \\ \hline
    Datos inválidos                    & Sistema muestra errores específicos por campo.                              \\ \hline
    Sin sesión o sin rol administrador & Sistema redirige a \texttt{/login} o muestra error de permisos.             \\ \hline
    Error del sistema                  & Mensaje genérico y registro en log.                                         \\ \hline
\end{longtable}

\subsubsection*{Flujo normal}
\begin{enumerate}[label=FN-\arabic*]
    \item Administrador accede a la sección “Gestión de usuarios”.
    \item Selecciona acción: consultar, modificar rol, desactivar o reactivar usuario.
    \item Sistema valida sesión activa y rol administrador.
    \item Administrador ingresa o selecciona datos del usuario (correo, nombre, rol).
    \item Sistema valida datos y persiste cambios en la BD.
    \item Sistema muestra confirmación y actualiza la lista de usuarios.
\end{enumerate}

\subsubsection*{Flujos alternativos}
\begin{enumerate}[label=FA-\arabic*]
    \item Usuario inexistente: sistema no encuentra el usuario en la BD; muestra error y no permite acción.
    \item Datos inválidos: sistema detecta campos vacíos o incorrectos; muestra errores específicos y mantiene formulario visible.
    \item Reactivar usuario: administrador selecciona usuario inactivo; sistema cambia estado a activo y lo muestra en la lista.
\end{enumerate}

\subsubsection*{Flujos excepcionales}
\begin{enumerate}[label=FE-\arabic*]
    \item Sin sesión o sin rol administrador: sistema detecta falta de permisos; redirige a \texttt{/login} o muestra error de autorización.
    \item Error en BD: sistema captura excepción, registra log y muestra mensaje genérico invitando a intentarlo de nuevo.
\end{enumerate}

% ============================
% CU14
% ============================

\subsection{CU14: Gestionar promociones y descuentos}
\subsubsection*{Objetivo}
Permitir al administrador dar de alta, modificar, desactivar o reactivar promociones y descuentos en el sistema, aplicándolos a productos o categorías según reglas de negocio.

\subsubsection*{Actores}
\begin{longtable}{|l|p{10cm}|}
    \hline
    Administrador & Gestiona las promociones y descuentos (alta, modificación, baja lógica, reactivación). \\ \hline
    Sistema       & Valida datos, persiste cambios en la BD y aplica reglas de negocio.                    \\ \hline
\end{longtable}

\subsubsection*{Condiciones iniciales}
\begin{itemize}[label=--]
    \item El administrador tiene sesión activa en el sistema.
    \item BD y servicio web disponibles.
    \item Existen productos o categorías registradas en la BD.
\end{itemize}

\subsubsection*{Condiciones de salida}
\begin{longtable}{|l|p{10cm}|}
    \hline
    Éxito                              & Promoción creada, modificada, desactivada o reactivada correctamente. \\ \hline
    Nombre duplicado                   & Sistema muestra error y no permite alta.                              \\ \hline
    Fechas inválidas                   & Sistema muestra error si la fecha de inicio es posterior a la de fin. \\ \hline
    Sin sesión o sin rol administrador & Sistema redirige a \texttt{/login} o muestra error de permisos.       \\ \hline
    Error del sistema                  & Mensaje genérico y registro en log.                                   \\ \hline
\end{longtable}

\subsubsection*{Flujo normal}
\begin{enumerate}[label=FN-\arabic*]
    \item Administrador accede a la sección “Gestión de promociones”.
    \item Selecciona acción: alta, modificación, desactivación o reactivación.
    \item Sistema valida sesión activa y rol administrador.
    \item Administrador ingresa o modifica datos de la promoción (nombre, tipo de descuento, porcentaje o monto, fecha de inicio y fin, productos o categorías aplicables).
    \item Sistema valida datos y persiste cambios en la BD.
    \item Sistema muestra confirmación y actualiza la lista de promociones.
\end{enumerate}

\subsubsection*{Flujos alternativos}
\begin{enumerate}[label=FA-\arabic*]
    \item Nombre duplicado: sistema detecta nombre repetido; muestra error y no permite alta.
    \item Fechas inválidas: sistema detecta que la fecha de inicio es posterior a la de fin; muestra error y no permite alta o modificación.
    \item Reactivar promoción: administrador selecciona promoción inactiva; sistema cambia estado a activo y la aplica nuevamente.
\end{enumerate}

\subsubsection*{Flujos excepcionales}
\begin{enumerate}[label=FE-\arabic*]
    \item Sin sesión o sin rol administrador: sistema detecta falta de permisos; redirige a \texttt{/login} o muestra error de autorización.
    \item Error en BD: sistema captura excepción, registra log y muestra mensaje genérico invitando a intentarlo de nuevo.
\end{enumerate}

% ============================
% TABLA DE TRAZABILIDAD
% ============================

\section{Tabla de trazabilidad}

\renewcommand{\arraystretch}{1.0}
\begin{tabularx}{\linewidth}{|p{2.5cm}|X|X|X|X|}
    \hline
    \textbf{CU}                              & \textbf{OE relacionado} & \textbf{RF relacionados} & \textbf{RN aplicables}  & \textbf{RNF aplicables} \\ \hline
    CU01: Registrar cuenta                   & OE1                     & RF01                     & RN1, RN2, RN4           & RNF02                   \\ \hline
    CU02: Iniciar sesión                     & OE1                     & RF02                     & RN3, RN4, RN5, RN6, RN8 & RNF02, RNF07            \\ \hline
    CU02b: Cerrar sesión                     & OE1                     & RF02b                    & RN4, RN5, RN6           & RNF02 
\\ \hline
    CU03: Ver catálogo                       & OE2                     & RF03a                    & RN9, RN13               & RNF03, RNF08            \\ \hline
    CU04: Buscar productos                   & OE2                     & RF03b                    & RN9, RN13               & RNF03, RNF08            \\ \hline
    CU05: Filtrar por categoría              & OE2                     & RF03c                    & RN10--RN13              & RNF03, RNF08            \\ \hline
    CU06: Ver detalle de producto            & OE2                     & RF03a, RF03b             & RN14--RN17, RN24--RN26b & RNF03, RNF08            \\ \hline
    CU07: Agregar producto al carrito        & OE3                     & RF04                     & RN18--RN20, RN22        & RNF07                   \\ \hline
    CU08: Gestionar carrito                  & OE3                     & RF04                     & RN18--RN20, RN22        & RNF07, RNF08            \\ \hline
    CU09: Realizar pago y confirmar orden    & OE4                     & RF05, RF06               & RN21--RN23, RN24--RN26b & RNF01, RNF07            \\ \hline
    CU10: Consultar historial de órdenes     & OE4                     & RF07                     & RN7, RN8                & RNF07, RNF08            \\ \hline
    CU11: Gestionar productos                & OE5                     & RF08, RF12               & RN14--RN17, RN27, RN28  & RNF02, RNF05            \\ \hline
    CU12: Gestionar categorías               & OE5                     & RF08                     & RN10--RN13, RN27, RN28  & RNF02, RNF05            \\ \hline
    CU13: Gestionar usuarios                 & OE1, OE5                & RF09                     & RN7, RN8, RN27, RN28    & RNF02, RNF05            \\ \hline
    CU14: Gestionar promociones y descuentos & OE5                     & RF11                     & RN24--RN26b, RN27, RN28 & RNF02, RNF05            \\ \hline
\end{tabularx}

\clearpage % fuerza salto de página si la imagen es muy grande

\section{Diagramas UML}

\subsection{Diagrama de Casos de Uso}
\begin{figure}[htbp]
    \centering
    % Escala la imagen al ancho del texto y ajusta altura si es necesario
    \includegraphics[width=\textwidth,height=0.7\textheight,keepaspectratio]{dcu.png}
    \caption{Diagrama de Casos de Uso del sistema Nozama}
    \label{fig:dcu}
\end{figure}

\clearpage % fuerza salto de página si la imagen es muy grande

\subsection{Diagrama de Clases}
\begin{figure}[htbp]
    \centering
    \includegraphics[width=\textwidth,height=0.7\textheight,keepaspectratio]{dc.png}
    \caption{Diagrama de Clases del sistema Nozama}
    \label{fig:dc}
\end{figure}

\end{document}

